% !TEX root = hw11-answers.tex
\chead{\textbf{Exercises week 11}}
\begin{document}
\flushleft\textbf{Topic: Causal Discovery}

\section{Reichenbach's Principle of Common Cause}
Reichenbach's Principle of Common Cause states that if two events are dependent,
then one must cause the other or the events are confounded (or any combination of these three possibilities).
We can make this precise for simple SCMs in the following way. Assume that $\C{M}$
is a simple SCM with two observed endogenous variables $X,Y$ (and possibly other
latent variables as well, but no exogenous input nodes). Prove that 
$X \not\Indep_{P_{\C{M}}} Y$
implies that $X \to Y$,
$X \ot Y$ or $X \oto Y$ in $G^{\{X,Y\}}(\C{M})$. (1 pt)

\begin{gradedsolution}
Assume that $X \to Y$, $X \ot Y$ and $X \oto Y$ \textbf{not} in $G^{\{X,Y\}}(\C{M})$. Then $X$ is d-separated from $Y$, and thus by Markov property $X \Indep_{P_{\C{M}}} Y$, which is a contradiction.
\end{gradedsolution}

\section{Randomized Controlled Trial assumptions}
Two key assumptions of {Randomized Controlled Trials with treatment variable $C$ and  are:
\begin{enumerate}
	\item outcome $X$ does not cause treatment $C$; (1/2 pt)
	\item outcome $X$ and treatment $C$ are not confounded, i.e., the values for
	 the treatment variable are assigned independently of other (latent) factors that may influence
	 the outcome. (1/2 pt)
  \end{enumerate}

Think like a conspiracy theorist and imagine situations in which the first assumption is not valid. Do the same for the second assumption.

\begin{gradedsolution}
	For example:
	\begin{enumerate}
		\item In a trial with a placebo, the doctor overseeing the coinflip might override the coinflip to assign patients to the non-placebo who are somehow more likely to benefit.
		\item When choosing between treatment $A$ and $B$, doctors might prefer treatment $B$ on less severe cases (see Kidney stone example of Simpons Paradox). 
	\end{enumerate}
\end{gradedsolution}

\section{Marginalization preserves separation}
By marginalizing out all variables in $V \cup W$ except $a$ and $b$, give short proofs of:
\begin{enumerate}
	\item Proposition 9.1.2: Let $\C{M} = \langle \Sin, \Send, \Srnd, \CX, P, f \rangle$ be a simple SCM with graph $G = G(\C{M})$.
	Let $a \in V \cup W \cup J$ and $b \in V$.
	If $a \notin \Anc^G(b)$, then:
	$$P_{\C{M}}(X_b \given \doit(X_{J\cup \{a\}})) = P_{\C{M}}(X_b \given \doit(X_{J\setminus \{a\}})).$$ (1 pt)
	\item Proposition 9.3.3:  Let $\C{M} = \langle \Sin, \Send, \Srnd, \CX, P, f \rangle$ be a simple SCM with graph $G = G(\C{M})$.
	Let $a \ne b \in V$.
	If $b \notin \Anc^G(a)$ and $a$ and $b$ are not confounded according to $\C{M}$, then
	%$b \Perp^\sigma_{G_{\doit(I_a)}} I_a \given \{a\} \cup J$.
	%Rule 2 of the do-calculus gives 
	  $P_{\C{M}}(X_b \given \doit(X_a), \doit(X_J)) = P_{\C{M}}(X_b \given X_a, \doit(X_J))$. (1 pt)
\end{enumerate}

\begin{gradedsolution}
\begin{enumerate}
	\item After marginalization we only retain variables $a$ and $b$ and all variables in $J$. Let us first consider the case $a \in V \cup W$. We will show that 
	$$b \Perp^\sigma_{G_{\doit(I_a)}} I_a \mid J \setminus \{a\}.$$
	Since $a$ is not an ancestor of $b$ and all edges from $J$ are outgoing, $a$ can only have incoming edges. Consider a walk from $b$ to $I_a \cup J$. Any walk from $b$ to $J$ is $\sigma$-blocked since we condition on the end points. Any walk from $b$ to $I_a$ must go through $a$, which is a collider since it only has incoming edges and is not conditioned on. Therefore the separation holds.
	Now consider the case $a \in J$. We will show that
	$$b \Perp^\sigma_{G} a \mid J \setminus \{a\}.$$
	Since $a \in J$, the only possible edge from $a$ is an edge to the only non-input variable $b$. But since $a \notin \Anc^G(b)$, $a$ has no edges. Consider a walk from $b$ to $J$. The node $a$ can not be an endpoint since it has no edges. On the other hand if the end point is in $J \setminus \{a\}$ it is $\sigma-$blocked by our conditioning. Thus the $\sigma$-separation holds.\\
	In both cases, rule 3 of do-calculus yields the result.
	\item After marginalization we only retain variables $a$ and $b$ and all variables in $J$. We will show that the assumptions imply $b \Perp^\sigma_{G_{\doit(I_a)}} I_a \given \{a\} \cup J$.
	
	Since $b \notin \Anc^G(a)$ and $a$ and $b$ are not confounded, the we have either $a \to b$ or no edge between $a$ and $b$. Consider a walk in $G_{\doit(I_a)}$ between $b$ and $I_a \cup J$. It cannot contain nodes from $J$ as such a node would either be an end node of the walk ($\sigma$-blocking it) or a non-collider node pointing only to nodes in another strongly connected component of $G_{\doit(I_a)}$
	($\sigma$-blocking the walk). Thus the only possible walk is $I_a \to a \to b$, which is $\sigma$-blocked by conditioning on $a$. Hence the $\sigma$-separation holds.

	Invoking rule 2 of the do-calculus then yields the result.
\end{enumerate}
\end{gradedsolution}

\section{Y-structures}
Prove proposition 9.8.1 by using the global Markov property and the faithfulness assumption, and showing that the only (cyclic or acyclic) graphs that are compatible with the observed
conditional independences are the ones in Figure 11, similarly as how is done in the proof of LCD (proposition  9.7.1). (1 pt)

\begin{gradedsolution}
	We have the conditional (in)dependencies:

	$$\begin{array}{lll}
		X_1 \not\Indep X_4,  & X_2 \not\Indep X_4, & X_1 \Indep X_2, \\
		X_1 \Indep X_4 \given X_3, & X_2 \Indep X_4 \given X_3, & X_1 \not\Indep X_2 \given X_3,
	  \end{array}$$
	
	By faithfulness assumption, from $X_1 \Indep X_2$, $X_1 \Indep X_4 \given X_3$ and $ X_2 \Indep X_4 \given X_3$ it respectively follows that there can be no direct edge between $X_1$ and $X_2$, between $X_1$ and $X_4$ and between $X_2$ and $X_4$, because otherwise the corresponding $\sigma$-separations do not hold. However, since $X_1 \not\Indep X_4$ and  $X_2 \not\Indep X_4$, there must exists a walk from $X_1$ to $X_4$ and a walk from $X_2$ to $X_4$, which by necessity must go through the only other node $X_3$. These facts together imply that there can only be edges between $X_3$ and the other nodes, and the graph is connected.

	Since $X_1 \Indep X_2$ and $X_1 \not\Indep X_2 \given X_3$, it follows that $X_3$ is a collider on the walk from $X_1$ to $X_2$. Therefore, any edges between $X_1$ and $X_3$ and $X_2$ and $X_3$ must have an arrowhead pointing into $X_3$.

	Finally, since $X_1 \not\Indep X_4$ and $X_1 \Indep X_4 \given X_3$, conditioning on $X_3$ must block the path from $X_1$ to $X_4$, therefore it cannot be a collider on this walk. As we have shown, edges between $X_1$ and $X_3$ must have an arrow pointing into $X_3$. Therefore, the edge between $X_4$ and $X_3$ cannot have an arrow pointing into $X_3$. It follows that $X_3 \to X_4$.
	
	With these facts taken together, it is easy to see that the only possible graphs are those drawn in figure 3.
	

\end{gradedsolution}

\end{document}
